\section{Introduction}
\label{sec:introduction}

The ultraviolet behavior of gravity remains a central open problem in theoretical physics.  Adding terms quadratic in curvature to the Einstein action gives a power-counting renormalizable theory and can yield asymptotic freedom at high energy~\cite{Stelle:1977ry,Avramidi:1985ki}.  The same terms introduce higher-derivative degrees of freedom, however, and their interpretation is inseparable from the question of negative-norm states.

Recent work has renewed interest in four-derivative field theories that admit a controlled perturbative expansion~\cite{Holdom:2024}.  In particular, the perfect-square (PS) scalar theory is related to the conformally flat limit of quadratic gravity.  In that limit, the local scale factor remains dynamical while graviton excitations decouple.  This gives a gravity-derived setting in which renormalization, indefinite metric structure, and proposed probability prescriptions can be studied without the full complexity of the spin-two sector~\cite{Anderson:2026,Bateman:2026}.

The aim of this paper is deliberately limited.  We do not claim a nonperturbative completion of full quadratic gravity.  Instead, we establish a numerical framework for asking whether the PS sector itself has a controlled continuum limit.  The primary questions are: (i) whether weak-coupling running agrees with perturbative expectations, (ii) whether correlation functions show stable finite-size scaling, and (iii) whether the Euclidean formulation admits a numerically well-defined continuum extrapolation.

Section~\ref{sec:model} specifies the continuum model and its relation to quadratic gravity.  Section~\ref{sec:numerics} presents the lattice and simulation strategy.  The results section is reserved for measured observables, uncertainty estimates, and continuum extrapolations.  Appendices record the discretization and observable definitions needed for reproduction.
